\documentclass[11pt]{article}
\usepackage[a4paper,margin=26mm]{geometry}
\usepackage[T1]{fontenc}
\usepackage{lmodern}
\usepackage{amsmath,amssymb,amsthm}
\usepackage{microtype}
\usepackage{array}
\usepackage{xurl}
\usepackage{hyperref}
\hypersetup{colorlinks=true,linkcolor=blue,urlcolor=blue,citecolor=blue,
 pdftitle={The integer group determinants of the Heisenberg group of order 125},
 pdfauthor={Horace Dong}}
\newtheorem{theorem}{Theorem}
\newtheorem{lemma}{Lemma}
\newtheorem{proposition}{Proposition}
\newtheorem{corollary}{Corollary}
\theoremstyle{definition}
\newtheorem{definition}{Definition}
\theoremstyle{remark}
\newtheorem{remark}{Remark}
\newcommand{\repoversion}{1.25.0}
\setlength{\emergencystretch}{1.5em}
\newcommand{\F}{\mathbb F_5}
\newcommand{\Z}{\mathbb Z}
\newcommand{\Tr}{\operatorname{Tr}}
\newcommand{\N}{\operatorname{N}}
\newcommand{\cS}{\mathcal S}
\newcommand{\vp}{v_\pi}

\title{The integer group determinants of the\\ Heisenberg group of order 125}
\author{Horace Dong\\ \small Independent researcher\\
\small\texttt{maxwelldhx@gmail.com}\\
\small ORCID: \href{https://orcid.org/0009-0002-7554-7548}{0009-0002-7554-7548}}
\date{22 November 2026}
\begin{document}
\maketitle

\begin{abstract}
Let $H_5$ be the Heisenberg group of order $125$, the group of upper
unitriangular $3\times3$ matrices over $\mathbb F_5$. We determine the set
$\cS(H_5)$ of all values of its group determinant $\det\bigl(f(gh^{-1})\bigr)_{g,h}$
when $f$ ranges over all integer-valued functions on $H_5$; this answers
Olga Taussky-Todd's integer group determinant question for this group.
Mossinghoff and Pinner had shown that the values prime to $5$ are the
integers $m$ with $m^4\equiv1\pmod{125}$, and that every nonzero value
divisible by $5$ is divisible by $5^{28}$. We show that every multiple of
$5^{33}$ is a value, and that a nonzero integer whose $5$-adic valuation lies
between $28$ and $32$ is a value if and only if it has a prime factor
$\ell\equiv1\pmod5$ that is \emph{visible}: a generator of a prime ideal of
$\mathbb Z[\zeta_5]$ above $\ell$ is not congruent modulo $5$ to a unit of
$\mathbb Z[\zeta_5]$. In particular, the least positive value divisible by $5$
is $11\cdot5^{28}=409781932830810546875$, attained by $1+x+y+y^2+xy^2$. The
classification and the minimum are formally verified in Lean~4. The proofs
were found and formalized by AI agents working under the author's direction;
their role is described at the end of the paper.
\end{abstract}

\section{Introduction}\label{sec:intro}

For a finite group $G=\{g_1,\dots,g_k\}$ and variables $a_g$ ($g\in G$), the
\emph{group determinant} is $\det\bigl(a_{g_ig_j^{-1}}\bigr)_{i,j}$, with $i$
indexing rows and $j$ columns~\cite[(1)]{mp}. At the meeting of the American
Mathematical Society in Hayward, California, in April 1977, Olga Taussky-Todd
asked which integers occur as circulant determinants with integer
entries~\cite[Section~1]{ps}; the same question can be asked for every finite group
$G$, namely to determine
\[
\cS(G)=\bigl\{\det\bigl(a_{g_ig_j^{-1}}\bigr)\ :\ a_g\in\Z\ \text{for all } g\in G\bigr\}
\]
\cite[(2)]{mp}. Complete answers are known for all groups of order at most
$14$~\cite{ps} and, according to~\cite{pan25}, for all groups of order less
than $20$ and both groups of order $25$. Writing an element of the group ring
as $F=\sum_g a_g g$, the matrix above is the matrix of left multiplication by
$F$ on $\Z[G]$ in the basis $G$, and we write $M(F)$ for its determinant.

Mossinghoff and Pinner~\cite{mp} studied the Heisenberg group
$H_p=\{x^ay^bz^c\}$ of order $p^3$, $p$ an odd prime, in which $z$ is central
and $yx=xyz$. Throughout we cite~\cite{mp} by the numbering of
arXiv:2108.04870v1, the only arXiv version; we could not access the journal
version (Michigan Math.\ J.\ 74 (2024), 551--569) to compare its numbering.
For $p=5$ their results say:
\begin{itemize}
\item the values prime to $5$ are exactly the integers $m$ with
$m^4\equiv1\pmod{125}$~\cite[Corollary~3.2]{mp};
\item if $5\mid M(F)$ then $5^{28}\mid M(F)$, and $5^{28}\,\|\,M(F)$ occurs
\cite[Theorem~3.4]{mp}.
\end{itemize}
They determined $\cS(H_3)$ completely~\cite[Theorem~3.6]{mp}; for $p\ge5$ the
multiples of $p$ in $\cS(H_p)$ were left undetermined. Their example for
$p=5$, $F=5-x+2y-y^2$~\cite[proof of Theorem~3.4, with $A=2$]{mp}, has
\[
M(F)=5^{28}\cdot7^{20}\cdot11\cdot101\cdot103^{20}\cdot131\cdot181\cdot191\cdot311,
\]
a cofactor of about $2.3\cdot10^{69}$ (our computation, by exact elimination
on the $125\times125$ matrix, confirmed by the final reviewer of the
minimum); so the
source settles the $5$-adic valuation $28$ but not the least value.

\paragraph{Visible primes.}
Let $\zeta$ be a primitive fifth root of unity, $K=\mathbb Q(\zeta)$,
$R=\Z[\zeta]$ its ring of integers, $\pi=1-\zeta$, and $\N=\N_{K/\mathbb Q}$
the norm, the product of the four conjugates. Let $U\subseteq(R/5R)^\times$ be
the image of the unit group $R^\times$.

\begin{definition}\label{def:visible}
A rational prime $\ell$ is \emph{visible} if $\ell\equiv1\pmod5$ and there is
$\lambda\in R$ with $\N(\lambda)=\ell$ whose residue class modulo $5R$ does not
lie in $U$.
\end{definition}

\noindent By Lemma~\ref{lem:choice}, ``there is'' can be replaced by ``for
every'': for a prime $\ell\equiv1\pmod 5$, either all elements of norm $\ell$
have residues in $U$ or none has. The visible primes below $600$ are
\[
11,31,41,61,71,101,131,151,181,191,241,251,271,311,331,401,431,491,541,571;
\]
the primes $\ell\equiv1\pmod5$ below $600$ that are not visible are
$211,281,421,461,521$. For example $\N(2+\zeta)=\Phi_5(-2)=11$ and the residue
of $2+\zeta$ is not in $U$.

\begin{theorem}\label{thm:main}
\begin{multline*}
\cS(H_5)=\{m\in\Z:\ m^4\equiv1\!\!\pmod{125}\}\ \cup\ 5^{33}\Z\\
\cup\ \{5^{28}n:\ n\in\Z\setminus\{0\},\ \text{some visible prime divides } n\}.
\end{multline*}
Equivalently, $0\in\cS(H_5)$, and for an integer $m\ne0$ with $v=v_5(m)$:
if $v=0$, then $m\in\cS(H_5)$ if and only if $m^4\equiv1\pmod{125}$; if
$1\le v\le27$, then $m\notin\cS(H_5)$; if $28\le v\le32$, then
$m\in\cS(H_5)$ if and only if $m$ has a visible prime factor; if $v\ge33$, then
$m\in\cS(H_5)$.
\end{theorem}

\begin{corollary}\label{cor:min}
The least positive element of $\cS(H_5)$ divisible by $5$ is
\[11\cdot5^{28}=409781932830810546875.\]
It is the group determinant of $1+x+y+y^2+xy^2$.
\end{corollary}

\begin{corollary}\label{cor:33}
$5^{32}\notin\cS(H_5)$. Hence $33$ is the least integer $k\ge0$ with
$5^k\Z\subseteq\cS(H_5)$.
\end{corollary}

For instance $11\cdot5^{28}$, $31\cdot5^{30}$ and $211\cdot11\cdot5^{28}$ are
values, while $211\cdot5^{28}$, $3\cdot5^{29}$ and $5^{28}$ are not.
Theorem~\ref{thm:main} and Corollary~\ref{cor:min} are formally verified in
Lean~4 (Section~\ref{sec:lean}); Corollary~\ref{cor:min} is a special case of
Theorem~\ref{thm:main}, but it was stated and verified separately, before
Theorem~\ref{thm:main}. Corollary~\ref{cor:33} follows from
Theorem~\ref{thm:main} in two lines (Section~\ref{sec:proofs}); it is not a
separate Lean statement. The second sentence of Theorem~\ref{thm:main} is an
immediate reformulation of the first (a visible prime is not $5$); only the
first is the Lean statement.

\begin{remark}\label{rem:cft}
Since $R$ is a principal ideal domain, the ray class group of $K$ modulo $5R$ is
$(R/5R)^\times/U\cong\Z/5$ (Lemma~\ref{lem:units}); so $\ell\equiv1\pmod5$ is
visible exactly when the primes above $\ell$ do not split completely in the
ray class field modulo $5R$, and by Chebotarev's theorem the visible primes
have density $4/5$ among the primes $\ell\equiv1\pmod5$ ($157$ of the $194$
below $6000$). As a purely numerical observation, not proved here: for these
$194$ primes (and, as a final reviewer checked, for all $563$ below $20{,}000$),
$\ell$ is visible if and only if $(1+r+r^{-1})^{(\ell-1)/5}\not\equiv1\pmod\ell$
for a primitive fifth root of unity $r$ modulo $\ell$, that is, if and only if
the golden-ratio unit $1+\zeta+\zeta^{-1}$ is not a fifth power modulo a prime
above $\ell$. Nothing in this remark is used below or stated in Lean.
\end{remark}

\paragraph{Prior work and what is new.}
The factorization of $M(F)$ used below, the congruence that controls the
five-dimensional factor, the description of the values prime to $5$ and the
bound $5^{28}$ are due to Mossinghoff and Pinner~\cite{mp}; the explicit values
prime to $5$ come from their Theorem~3.3. The constructions start from the three
group ring elements of the abelian group $C_5\times C_5$ found by
Panraksa~\cite[Proposition~3.1]{pan25}, which he used to prove
$\cS(C_5^2)\cap5\Z=5^8\Z$. For $C_7\times C_7$, Panraksa~\cite{pan49}
described the values at the two lowest valuations by conditions on ideals of
$\Z[\zeta_7]$, which is a precedent for a criterion of the kind found here; for
$C_p\times C_p$ with $p\ge7$ see~\cite{panp}; we borrow the words visible and
invisible from~\cite{pan49}, where they are used for the analogous condition
on primes of $\Z[\zeta_7]$. The lattice isomorphism algorithms of~\cite{bhj} decide membership in
$\cS(G)$ in principle (an observation made during this work, not used here),
but give no explicit description. What is new
here, as far as the searches described at the end found, is the description of
the multiples of $5$ in $\cS(H_5)$: the inclusion $5^{33}\Z\subseteq\cS(H_5)$,
the visible-prime criterion at the valuations $28$ to $32$, and the minimum of
Corollary~\ref{cor:min}.

\section{The group determinant of \texorpdfstring{$H_5$}{H5}}\label{sec:factor}

\paragraph{The group.} We write the elements of $H_5$ as $x^ay^bz^c$ with
$a,b,c\in\F$; from $yx=xyz$ and the centrality of $z$,
\[
x^ay^bz^c\cdot x^{a'}y^{b'}z^{c'}=x^{a+a'}y^{b+b'}z^{c+c'+ba'} .
\]
In the matrix model of~\cite[p.~4]{mp}, $x^ay^bz^c$ is the unitriangular
matrix with entries $b$, $a$, $c$ in positions $(1,2)$, $(2,3)$, $(1,3)$. A
group ring element is $F=\sum f(a,b,c)\,x^ay^bz^c$ with $f:H_5\to\Z$
arbitrary, and $M(F)=\det\bigl(f(gh^{-1})\bigr)_{g,h\in H_5}$. Its
\emph{abelianization} is
\[
P(x,y)=\sum_{a,b\in\F}p_{ab}\,x^ay^b,\qquad p_{ab}=\sum_{c\in\F}f(a,b,c),
\]
an element of $\Z[x,y]/(x^5-1,y^5-1)$, and $t=P(1,1)=\sum_gf(g)$ is the
augmentation. We write $P_X=\sum a\,p_{ab}$ and $P_Y=\sum b\,p_{ab}$, read
modulo $5$.

\paragraph{The factorization.}
By Frobenius' theorem, $M(F)=\prod_\rho\det\rho(F)^{\deg\rho}$ over the
irreducible complex representations $\rho$ of $H_5$~\cite[\S3.1]{mp}. There
are $25$ linear characters, $x^ay^bz^c\mapsto\zeta^{ia+jb}$, with values
$P(\zeta^i,\zeta^j)$, and four representations $\phi_\lambda$ of degree $5$,
$\lambda=\zeta^k$ ($1\le k\le4$), in which $x$ acts by the cyclic shift
$S e_j=e_{j+1}$, $y$ by $T=\mathrm{diag}(1,\lambda,\dots,\lambda^4)$ and $z$
by $\lambda$. Thus $\phi_\zeta(F)=\sum_{a,b}r_{ab}S^aT^b$ with
$r_{ab}=\sum_cf(a,b,c)\zeta^c$, and we put $D=\det\phi_\zeta(F)\in R$. The
other three representations give the Galois conjugates of $D$, and
Mossinghoff and Pinner obtain $M(F)=M_1M_2^5$ with
$M_1=\prod_{i,j}P(\zeta^i,\zeta^j)$ and $M_2=\N(D)$. Grouping the $24$
nontrivial characters into the six orbits of $\F^\times$, each of which is a
full set of Galois conjugates, and choosing the representatives $(1,0)$ and
$(j,1)$, $j\in\F$, we get
\begin{equation}\label{eq:factor}
M(F)=t\cdot\N(\gamma)\cdot\prod_{j\in\F}\N(\alpha_j)\cdot\N(D)^5,
\qquad \gamma=P(\zeta,1),\quad \alpha_j=P(\zeta^j,\zeta).
\end{equation}
All factors are integers, and norms of elements of $R$ are $\ge0$. We call
$\gamma$ and the $\alpha_j$ the values on the six \emph{rays}. The Lean files
prove~\eqref{eq:factor} from the matrix itself (Section~\ref{sec:lean}).

\paragraph{Arithmetic of $R$.} We use the following standard facts. (a)~$5R=\pi^4R$,
$R/\pi R=\F$, $\N(\pi)=5$, and $v_5(\N(\alpha))=\vp(\alpha)$ for
$\alpha\in R$; moreover $\zeta^n\equiv1-n\pi+\binom n2\pi^2\pmod{\pi^3}$.
(b)~$\Tr(\pi R)\subseteq5\Z$, since $\Tr(\zeta^i-\zeta^{i+1})\in\{0,\pm5\}$;
hence $\Tr(\pi^5R)=\Tr(5\pi R)\subseteq25\Z$ and
$\Tr(\pi^9R)\subseteq125\Z$; also $\Tr(\pi^4)=5$ and $\Tr(\pi^8)=-275$.
(c)~$R$ is a principal ideal domain (the class number of $K$ is one). (d)~Every
unit $u\in R^\times$ has $\N(u)=1$, and $\bar u/u$ is a root of unity (all its
conjugates have absolute value~$1$); comparing residues modulo $\pi$ gives
$\bar u=\zeta^ku$ for some $k$. Let $\varphi=1+\zeta+\zeta^{-1}=-\zeta^2-\zeta^3$,
the golden-ratio unit; $\varphi\equiv3+\pi^2+\pi^3\pmod5$.

\begin{lemma}\label{lem:units}
$(R/5R)^\times$ has $500$ elements. Every element of it can be written
uniquely as $v\,(1+d\pi^3)$ with $v$ in the subgroup generated by $\zeta$ and
$\varphi$ and $d\in\F$, and $d$ is a surjective homomorphism
$(R/5R)^\times\to\F$ whose kernel is $U$; thus $|U|=100$. For the
automorphism $\sigma_k:\zeta\mapsto\zeta^k$ of $K$, $d(\sigma_k\lambda)=k^3d(\lambda)$;
in particular complex conjugation sends $d$ to $-d$.
\end{lemma}

\begin{proof}
$R/5R=\F[\pi]/(\pi^4)$. Multiplying by a power of $\varphi\equiv3\pmod\pi$
makes the constant term $1$; by a power of $\zeta\equiv1-\pi$ removes the term
in $\pi$; by a power of $\varphi^4\equiv1+3\pi^2+3\pi^3\pmod 5$ removes the
term in $\pi^2$. This leaves $1+d\pi^3$, and
$(1+d\pi^3)(1+d'\pi^3)\equiv1+(d+d')\pi^3$. For uniqueness and for $U$ it
suffices to show that a unit $u\equiv1+d\pi^3\pmod5$ has $d=0$. By~(d),
$\bar u=\zeta^ku$; since $\bar\pi\equiv-\pi\pmod{\pi^2}$ we have
$\bar u\equiv1-d\pi^3$, while $\zeta^ku\equiv1-k\pi\pmod{\pi^2}$; so $k=0$ and
then $1-d\pi^3\equiv1+d\pi^3$, i.e.\ $d=0$. Finally
$\sigma_k(\pi)=1-\zeta^k\equiv k\pi\pmod{\pi^2}$, so
$\sigma_k(1+d\pi^3)\equiv1+k^3d\pi^3$, and $\sigma_k$ maps units to units.
\end{proof}

\begin{lemma}[choice independence]\label{lem:choice}
Let $\ell\equiv1\pmod5$ be prime. If $\lambda,\lambda'\in R$ both have norm
$\ell$, then $d(\lambda)=0$ if and only if $d(\lambda')=0$.
\end{lemma}

\begin{proof}
An element of prime norm generates a prime ideal of norm $\ell$. The prime
ideals above $\ell$ are permuted transitively by the Galois group, and two
generators of the same ideal differ by a unit, so $\lambda'=u\,\sigma_k(\lambda)$
with $u\in R^\times$, and $d(\lambda')=k^3d(\lambda)$ by
Lemma~\ref{lem:units}.
\end{proof}

\begin{lemma}[prime extraction]\label{lem:extract}
Let $\mu\in R$ with $5\nmid\N(\mu)$ and $\mu\bmod 5\notin U$. Then some
visible prime divides $\N(\mu)$.
\end{lemma}

\begin{proof}
Factor $\mu$ into prime elements (by~(c)). Since $d$ is a homomorphism, some
prime factor $\lambda$ has $d(\lambda)\ne0$. If $\bar\lambda=u\lambda$ with
$u$ a unit, then $-d(\lambda)=d(\bar\lambda)=d(\lambda)$, so $d(\lambda)=0$, a
contradiction. Hence $\lambda$ and $\bar\lambda=\sigma_2^2\lambda$ are not
associate, and then neither are any two of the four conjugates of $\lambda$
(the Galois group is cyclic, generated by $\sigma_2$). All four divide the
rational prime $\ell$ below $\lambda$, so their product $\N(\lambda)$ divides
$\ell$ and $\N(\lambda)=\ell$. As $\N(\lambda)\equiv(\lambda\bmod\pi)^4\equiv1\pmod5$,
$\ell$ is visible, and $\ell\mid\N(\mu)$.
\end{proof}

\section{The \texorpdfstring{$5$}{5}-adic floor and the five-dimensional factor}\label{sec:floor}

\begin{lemma}[normalization]\label{lem:normalize}
For every $f$ there is $f'$ with $M(f')=M(f)$ whose abelianization has
$P_X\equiv0\pmod5$.
\end{lemma}

\begin{proof}
Every automorphism of $H_5$ permutes the basis of $\Z[H_5]$ and so preserves
$M$. Any basis of $\F^2=H_5/\langle z\rangle$ lifts to an automorphism, since
lifts of two independent elements have order $5$ and their commutator is a
nonzero power of $z$. The automorphism acts on the vector $(P_X,P_Y)\in\F^2$ through the induced
linear map, and $\mathrm{GL}_2(\F)$ is transitive on nonzero vectors, so we can
make $P_X\equiv0$.
\end{proof}

\noindent From now on in this section $P_X\equiv0$. By~(a), for a ray value,
\begin{equation}\label{eq:taylor}
P(\zeta^u,\zeta^v)\equiv t-(uP_X+vP_Y)\,\pi+\sum_{a,b}p_{ab}\tbinom{ua+vb}{2}\pi^2\pmod{\pi^3}.
\end{equation}

\begin{lemma}[{the congruence of~\cite{mp}}]\label{lem:c0}
Write $\prod_{x^5=1}P(x,y)\equiv c_0+c_1y+\dots+c_4y^4\pmod{y^5-1}$, with
$c_i\in\Z$. Then $5\mid c_i$ for $i\ne0$,
\[
D\equiv c_0\pmod{\pi^5},\qquad 5c_0=t\,\N(\gamma)+\Tr\beta,\quad
\beta=\prod_{j\in\F}\alpha_j .
\]
\end{lemma}

\begin{proof}
The first two assertions are~(18) and the congruence after (15) and (16) in the
proof of~\cite[Theorem~3.1]{mp}, which holds before that proof splits into the
cases $p\mid c_0$ and $p\nmid c_0$; the Lean files prove both again, by
expanding $D$ along column selections. For the identity,
$5c_0=\sum_{y^5=1}\prod_{x^5=1}P(x,y)$. The term $y=1$ is
$P(1,1)\prod_{u\ne0}P(\zeta^u,1)=t\,\N(\gamma)$, and for $y=\zeta^j$, $j\ne0$,
the term is $\prod_uP(\zeta^u,\zeta^j)=\sigma_j(\beta)$.
\end{proof}

\begin{proposition}[{after~\cite[Theorem~3.4]{mp}}]\label{prop:floor}
Suppose $5\mid M=M(F)$, with $P_X\equiv0$. Then $5\mid t$, $\pi^2\mid\gamma$,
$\pi\mid\alpha_j$ for all $j$, $\pi^4\mid D$, and so $5^{28}\mid M$. If
moreover $M\ne0$ and $v_5(M)\le32$, then $P_Y\not\equiv0$, $\vp(\alpha_j)=1$
for all $j$, $\vp(D)=4$, and $125\nmid\Tr\beta$.
\end{proposition}

\begin{proof}
Modulo $\pi$ every character value is $\equiv t$, and $\phi_\zeta(F)\equiv
\sum p_{ab}S^a$, a circulant matrix with determinant
$\prod_{x^5=1}P(x,1)\equiv t^5$; so $M\equiv t^{125}\pmod\pi$,
and $5\mid M$ forces $\pi\mid t$, i.e.\ $5\mid t$. Then $t\equiv0\pmod{\pi^4}$,
and~\eqref{eq:taylor} gives $\gamma\equiv-P_X\pi\equiv0$ and
$\alpha_j\equiv-P_Y\pi\pmod{\pi^2}$. So $25\mid\N(\gamma)$ and
$125\mid t\,\N(\gamma)$, while $\beta\in\pi^5R$ gives $\Tr\beta\in25\Z$. By
Lemma~\ref{lem:c0}, $25\mid5c_0$, so $\pi^4\mid c_0$ and $\pi^4\mid D$.
Counting valuations in~\eqref{eq:factor}: $1+2+5\cdot1+5\cdot4=28$.

Now let $M\ne0$ and $v_5(M)\le32$. If also $P_Y\equiv0$, all six ray values
are divisible by $\pi^2$ and $v_5(M)\ge1+12+20=33$; so $P_Y\not\equiv0$ and
$\vp(\alpha_j)=1$. If $\pi^5\mid D$, then $v_5(M)\ge8+25=33$; so $\vp(D)=4$.
By Lemma~\ref{lem:c0}, $\pi^5\nmid c_0$, i.e.\ $25\nmid c_0$, so
$125\nmid5c_0$, and as $125\mid t\N(\gamma)$, $125\nmid\Tr\beta$.
\end{proof}

\begin{lemma}[{values prime to $5$; cf.~\cite[Theorem~3.1]{mp}}]\label{lem:coprime}
$M(F)\equiv c_0^{25}\pmod{125}$. Hence, if $5\nmid M(F)$, then
$M(F)^4\equiv1\pmod{125}$.
\end{lemma}

\begin{proof}[Proof (sketch; this is the route of \texttt{Coprime.determinant\_congruence})]
By Lemma~\ref{lem:c0}, $D=c_0+5w$ with $\pi\mid w$, so
$\N(D)\equiv c_0^4+5c_0^3\Tr w\equiv c_0^4\pmod{25}$ and
$\N(D)^5\equiv c_0^{20}\pmod{125}$. Put $C(y)=\sum c_iy^i$ and
$s=\sum_{i\ne0}c_i/5$; then $M_1=C(1)\,\N(C(\zeta))$ with $C(1)=c_0+5s$ and
$C(\zeta)=c_0+5w'$, where $\Tr w'=-s$ and $w'\equiv s\pmod\pi$. Expanding the
norm to second order and using $\Tr(\mu)\equiv-(\mu\bmod\pi)\pmod5$ gives
$M_1\equiv c_0^5\pmod{125}$. If $5\nmid M$, then $5\nmid c_0$ and
$M^4\equiv c_0^{100}\equiv1\pmod{125}$.
\end{proof}

\section{Necessity at the valuations 28 to 32: phases}\label{sec:necessity}

Assume $5\mid t$, $P_X\equiv0$ and $P_Y\not\equiv0$. By~\eqref{eq:taylor},
\begin{equation}\label{eq:AB}
\alpha_j\equiv A\pi+B(j)\pi^2\pmod{\pi^3},\qquad
A=-P_Y,\quad B(j)=\sum_{a,b}p_{ab}\tbinom{ja+b}{2},
\end{equation}
where $A\in\F^\times$ and $B$ is a polynomial function of $j\in\F$ of degree at
most $2$ (the terms $t$ and $jP_X\pi$ lie in $5R\subseteq\pi^3R$).

\begin{lemma}[phases]\label{lem:phase}
Let $u_j\in R^\times$ ($j\in\F$) with $\pi u_j\equiv A\pi+B(j)\pi^2\pmod{\pi^3}$,
where $A\ne0$ and $B$ has degree at most $2$. Then
$\overline{\pi u_j}=-\zeta^{k_j}\pi u_j$ with $k_j=2B(j)/A-1$, $\sum_jk_j=0$,
and $\Tr\prod_j\pi u_j=0$.
\end{lemma}

\begin{proof}
By~(d), $\bar u_j=\zeta^{r}u_j$, and $\bar\pi=-\zeta^{-1}\pi$; so
$\overline{\pi u_j}=-\zeta^{k_j}\pi u_j$ for some $k_j$. Write
$\alpha=\pi u_j=A\pi+B\pi^2+\dots$. Since
$\bar\pi\equiv-\pi-\pi^2\pmod{\pi^3}$, we get
$\bar\alpha\equiv-A\pi+(B-A)\pi^2$, while
$-\zeta^{k}\alpha\equiv-A\pi+(kA-B)\pi^2$; hence $k_jA=2B(j)-A$. Since
$\sum_{j\in\F}j^e=0$ in $\F$ for $e=0,1,2$, $\sum_jk_j=0$. So
$\beta'=\prod_j\pi u_j$ satisfies $\bar\beta'=-\zeta^{\sum k_j}\beta'=-\beta'$,
and as complex conjugation preserves the trace, $\Tr\beta'=0$.
\end{proof}

\begin{proposition}\label{prop:necessity}
If $5\mid M(F)$, then $5^{33}\mid M(F)$, or $M(F)=5^{28}n$ with $n\ne0$ and
some visible prime divides $n$.
\end{proposition}

\begin{proof}
By Lemma~\ref{lem:normalize} we may assume $P_X\equiv0$. Suppose $M\ne0$ and
$v_5(M)\le32$. By Proposition~\ref{prop:floor}, $5^{28}\mid M$,
$P_Y\not\equiv0$, $\alpha_j=\pi q_j$ with $\pi\nmid q_j$, and
$125\nmid\Tr\beta$. If every $q_j$ were congruent modulo $5$ to a unit $u_j$,
then $\alpha_j\equiv\pi u_j\pmod{5\pi}$, where $5\pi R=\pi^5R$; so $\pi u_j$
has the same expansion~\eqref{eq:AB} and Lemma~\ref{lem:phase} gives
$\Tr\prod\pi u_j=0$. Replacing the five factors one at a time changes the
product by elements of $\pi^{5+4}R$, so
$\Tr\beta\equiv\Tr\prod\pi u_j=0\pmod{125}$ by~(b), a contradiction. So
some $q_j$ has residue outside $U$; as $5\nmid\N(q_j)$,
Lemma~\ref{lem:extract} gives a visible prime $\ell\mid\N(q_j)$, and
$\N(q_j)$ divides $M$ by~\eqref{eq:factor}. Since $\ell\ne5$, $\ell$ divides
$n=M/5^{28}$.
\end{proof}

\noindent This is where the visible primes come from: when the residues of the
$q_j$ come from global units, the phases $k_j$ sum to zero and $D$ acquires an
extra factor $\pi$; a visible prime contributes a term $d\pi^3$, $d\ne0$, that
breaks this balance, and the constructions below use exactly this term.

\section{Sufficiency: seeds, one ray, and an exact lift}\label{sec:sufficiency}

\paragraph{The seeds.} Panraksa's elements of $\Z[C_5^2]$~\cite[(12)--(14),
Proposition~3.1]{pan25} are
\begin{gather*}
F_1=1+y+y^2+x+x^4y,\qquad F_5=\Phi_5(x)\Phi_5(y)-1-y+x+x^4,\\
F_2=1+y+y^2+y^3+x+xy+xy^2+x^4+x^4y+x^4y^4 .
\end{gather*}
Their augmentations are $5,10,25$; their value on the ray $(1,0)$ has norm
$25$ and the other five ray values have norm $5$~\cite[proof of
Proposition~3.1]{pan25}. Together with $-F_1(x,y^{-1})$ and $-F_2(x,y^{-1})$
this gives five elements $S_s$ with augmentations $5a_s$,
$a_s\in\{1,-1,2,-2,5\}$, a complete set of residues modulo $5$, and the same
ray norms. Exchanging the variables, put $P_s(x,y)=S_s(y,x)$. Then
$b_s:=P_s(\zeta,1)$ has norm $5$, and $\prod_{j\in\F}\N(P_s(\zeta^j,\zeta))=25\cdot5^4=5^6$.

\paragraph{Row shifts.} For $q:\F\to\Z$ put $Q(\zeta)=\sum_aq(a)\zeta^a$. Adding
$q(a)$ to every coefficient $p_{ab}$, $b\in\F$, of an abelianization adds
$q(a)x^a\Phi_5(y)$; this adds $5\sum_aq(a)$ to $t$ and $5Q(\zeta)$ to
$P(\zeta,1)$, and leaves every $P(\zeta^j,\zeta)$ unchanged, because
$\Phi_5(\zeta)=0$. So by~\eqref{eq:factor}, any $f$ whose abelianization is
$P_s(a,b)+q(a)$ has
\begin{equation}\label{eq:rows}
M(f)=5\Bigl(a_s+\sum_aq(a)\Bigr)\cdot\N\bigl(b_s+5Q(\zeta)\bigr)\cdot5^6\cdot\N(D)^5 .
\end{equation}

\begin{lemma}[exact lift]\label{lem:lift}
For every $s$ and every $q:\F\to\Z$ there is $f:H_5\to\Z$ with abelianization
$P_s(a,b)+q(a)$ and $D=5\tau$, where $\tau=\pi$ if $5\mid\sum_aq(a)$, and
otherwise $\tau$ is a unit with $\tau\equiv\sum_aq(a)\pmod\pi$. Thus
$\N(D)=5^5$ or $5^4$, respectively.
\end{lemma}

\begin{proof}
Fixing the abelianization fixes $\phi_\zeta(F)=\sum r_{ab}S^aT^b$ only up to
$r_{ab}\in p_{ab}+\pi R$, and every such choice is realized, because
$\{\sum_c f_c\zeta^c:\sum_cf_c=p\}=p+\pi R$. For each $s$ the Lean files
contain an explicit base lift $f_s$ ($125$ integers) with abelianization $P_s$
whose matrix $C_s=\phi_\zeta(f_s)$ has determinant $0$, together with the
cofactors $l_i$ of its column $0$ and elements $e_i\in R$ such that
$l_i\equiv1\pmod\pi$ and $\sum_il_ie_i=1$. Adding $g(a,c)$ to every
$f(a,b,c)$, $b\in\F$, changes only column $0$ of the matrix, because
$\sum_bT^b=\mathrm{diag}(5,0,0,0,0)$: its entry in row $a$ grows by
$5\sum_cg(a,c)\zeta^c$, and $p_{ab}$ grows by $\sum_cg(a,c)$. Let
$w=\sum_il_iq(i)\equiv\sum_iq(i)\pmod\pi$, choose $\tau$ as stated (so that
$\tau\equiv w\pmod\pi$; a power of $\varphi\equiv3$ gives every nonzero
residue), and put $x_i=(\tau-w)e_i\in\pi R$. Choose $g(i,\cdot)$ with
$\sum_cg(i,c)=q(i)$ and $\sum_cg(i,c)\zeta^c=q(i)+x_i$, which is possible
since $x_i\in\pi R$. Since the determinant is linear in column $0$,
$D=\det C_s+5\sum_il_i(q(i)+x_i)=5w+5(\tau-w)=5\tau$.
\end{proof}

\noindent As Lemma~\ref{lem:c0} requires, the residue of $D/5$ modulo $\pi$ is
determined by the abelianization (here it is $\sum_aq(a)\bmod5$).

\begin{proposition}[the tail]\label{prop:tail}
$5^{33}\Z\subseteq\cS(H_5)$.
\end{proposition}

\begin{proof}
Given $m\in\Z$, choose $s$ with $a_s\equiv m\pmod5$ and let $q\equiv k$ be
constant with $k=(m-a_s)/5$. Then $\sum q=5k$, $Q(\zeta)=k\Phi_5(\zeta)=0$,
and Lemma~\ref{lem:lift} gives $\N(D)=5^5$. By~\eqref{eq:rows},
$M(f)=5m\cdot5\cdot5^6\cdot5^{25}=5^{33}m$.
\end{proof}

\begin{proposition}[the layers $28$ to $32$]\label{prop:low}
If $\ell$ is visible and $\ell\mid n$, then $5^{28}n\in\cS(H_5)$.
\end{proposition}

\begin{proof}
Write $n=\ell k$. By Definition~\ref{def:visible} and Lemma~\ref{lem:units},
after multiplying by a unit there is $\eta\in R$ with $\N(\eta)=\ell$ and
$\eta=1+d\pi^3+5w$, $d\ne0$, $w\in R$. Choose $s$ with $a_s\equiv k+d\pmod 5$.
A direct computation shows $\pi^3b_s\in5R$ and
$\pi^3b_s/5\equiv-1\pmod\pi$ for all five seeds. Put
$W=d\,\pi^3b_s/5+wb_s$, so that $b_s+5W=\eta\,b_s$ and $W\equiv-d\pmod\pi$.
Every $W\in R$ equals $Q(\zeta)$ for integers $q(a)$ whose sum is any
prescribed integer $\equiv W\pmod\pi$ (add a constant to all $q(a)$); take
$\sum_aq(a)=k-a_s$, which is $\equiv-d\not\equiv0$ (this is where the class
$d$ of the visible prime is used). By Lemma~\ref{lem:lift},
$\N(D)=5^4$, and by~\eqref{eq:rows},
$M(f)=5k\cdot\N(\eta)\N(b_s)\cdot5^6\cdot5^{20}=5^{28}\ell k$.
\end{proof}

\begin{proposition}[{values prime to $5$; \cite[Theorem~3.3]{mp}}]\label{prop:coprime}
Every $m$ with $m^4\equiv1\pmod{125}$ lies in $\cS(H_5)$.
\end{proposition}

\begin{proof}
Adding $k\sum_gg$ to $F$ changes only the trivial character, so it changes
$M(F)$ by $125k$ times the product of the other factors. Following
Mossinghoff and Pinner's construction (which gives $a^{25}+125k$ for
$5\nmid a$; note $2^{25}\equiv57\pmod{125}$), the Lean files use two elements
with $M=1$ and $M=57$ whose nontrivial factors all equal $1$, so that the
shifts give $1+125\Z$ and $57+125\Z$; as $M(-F)=-M(F)$, this covers the four
solutions $\pm1,\pm57$ of $m^4\equiv1\pmod{125}$.
\end{proof}

\section{Proofs of the theorem and the corollaries}\label{sec:proofs}

\begin{proof}[Proof of Theorem~\ref{thm:main}]
Let $m=M(F)$. If $5\nmid m$, Lemma~\ref{lem:coprime} gives
$m^4\equiv1\pmod{125}$; if $5\mid m$, Proposition~\ref{prop:necessity} puts
$m$ in the second or third set. Conversely, the three sets lie in
$\cS(H_5)$ by Propositions~\ref{prop:coprime}, \ref{prop:tail}
and~\ref{prop:low}. The reformulation by valuations is immediate, since a
visible prime is not $5$.
\end{proof}

\begin{proof}[Proof of Corollary~\ref{cor:min}]
A positive multiple of $5$ in $\cS(H_5)$ is either a multiple of $5^{33}$ or
$5^{28}n$ with $n\ge1$ divisible by a visible prime; $11$ is the least prime
$\equiv1\pmod5$, and it is visible. So the minimum is
at least $11\cdot5^{28}$, and $11\cdot5^{28}<5^{33}$. For
$F=1+x+y+y^2+xy^2=f_0(y)+xf_1(y)$ with $f_0=1+y+y^2$, $f_1=1+y^2$,
formulas~(9) and~(10) of~\cite{mp} give
$D=\prod_{y^5=1}f_0(y)+\prod_{y^5=1}f_1(y)=3+2=5$, so $\N(D)=5^4$, and
$M_1=\prod_{y^5=1}\bigl(f_0(y)^5+f_1(y)^5\bigr)$. The factor $y=1$ is
$3^5+2^5=275$; for $y=\zeta$, $f_0^5+f_1^5=-5(1+2u)$ with
$u=\zeta^2+\zeta^3$, and $(1+2u)^2=5$ because $u^2+u-1=0$, so the four
conjugates multiply to $5^4\cdot25=5^6$. Hence $M_1=11\cdot5^8$ and
$M(F)=11\cdot5^8\cdot5^{20}=11\cdot5^{28}$. In terms of~\eqref{eq:factor},
$t=5$, $\N(P(\zeta,1))=\N(3+2\zeta)=55$, the exceptional ray $(0,1)$ has norm
$25$ and the other four have norm~$5$.
\end{proof}

\begin{proof}[Proof of Corollary~\ref{cor:33}]
$5^{32}$ is divisible by $5$, not by $5^{33}$, and $5^{32}/5^{28}=625$ has no
prime factor other than $5$; so $5^{32}\notin\cS(H_5)$ by
Theorem~\ref{thm:main}, and hence $5^k\Z\not\subseteq\cS(H_5)$ for $k\le32$,
while $5^{33}\Z\subseteq\cS(H_5)$. This deduction is not a separate Lean
statement.
\end{proof}

{\sloppy
\section{Formal verification}\label{sec:lean}

Theorem~\ref{thm:main} and Corollary~\ref{cor:min} are formally verified in
Lean~4 (version 4.34.1) with Mathlib~\cite{mathlib} (commit
\texttt{d13f23b7}). The $63$ files are in version~\repoversion{} of the
repository~\cite{repo}: the $37$ files of the minimum in
\texttt{lean/Research/TaskU/H5/} and the $26$ further files of the
classification in \texttt{lean/Research/TaskU/H5B/}; records of the freezing
and of the checks are in the folder
\texttt{certificates/heisenberg-125-determinants}.

\paragraph{The statements.} The statement file \texttt{Challenge.lean} of the
minimum, in ASCII transliteration (the \texttt{import} and docstring lines, the outer
\texttt{namespace TaskU.H5} lines, four auxiliary \texttt{private} lemmas about
the coordinates, and the proofs of the group axioms are omitted; the file uses
Unicode, e.g.\ $\langle\ \rangle$, $\mathbb Z$, $\exists$, $\wedge$, $\mid$;
\texttt{\^{}-1} stands for ${}^{-1}$ and \texttt{->} for $\to$):
{\footnotesize
\begin{verbatim}
structure H where
  a : ZMod 5
  b : ZMod 5
  c : ZMod 5
  deriving DecidableEq, Fintype
instance : Mul H where
  mul g h := <g.a + h.a, g.b + h.b, g.c + h.c + g.b * h.a>
instance : One H where
  one := <0, 0, 0>
instance : Inv H where
  inv g := <-g.a, -g.b, -g.c + g.b * g.a>
instance : Group H where ...          -- proofs of the group axioms omitted
def x : H := <1, 0, 0>
def y : H := <0, 1, 0>
def witness (g : H) : Int :=
  (if g = 1 then 1 else 0) + (if g = x then 1 else 0) +
  (if g = y then 1 else 0) + (if g = y ^ 2 then 1 else 0) +
  (if g = x * y ^ 2 then 1 else 0)
noncomputable def groupDet (f : H -> Int) : Int :=
  Matrix.det (Matrix.of fun g h : H => f (g * h^-1))
namespace Challenge
def witnessValue : Prop := groupDet witness = 11 * 5 ^ 28
def leastPositiveMultiple : Prop :=
  IsLeast {m : Int | 0 < m /\ 5 | m /\ exists f : H -> Int, groupDet f = m}
    (11 * 5 ^ 28)
end Challenge
\end{verbatim}}
\noindent The element $\langle a,b,c\rangle$ is $x^ay^bz^c$, and the
multiplication is the one displayed in Section~\ref{sec:factor}.
\texttt{groupDet f} is $M(F)$ with row $g$ and column $h$, as in~\cite[(1)]{mp};
\texttt{f} is an arbitrary function $H\to\Z$, one integer per group element,
negative values included. \texttt{IsLeast S v} means $v\in S$ and $v\le s$ for
all $s\in S$. So \texttt{witnessValue} is the value in Corollary~\ref{cor:min}
and \texttt{leastPositiveMultiple} is the minimality. With
\texttt{fullSpectrum} below, the frozen targets are
\texttt{TaskU.H5.Challenge.witnessValue},
\texttt{TaskU.H5.Challenge.leastPositiveMultiple} and
\texttt{TaskU.H5.ChallengeB.fullSpectrum}.

The statement file \texttt{ChallengeB.lean} of the classification (imports,
namespace and docstrings omitted; \texttt{==} stands for $\equiv$,
\texttt{[ZMOD n]} for congruence modulo $n$, \texttt{\%} for the remainder,
\texttt{!=} for $\ne$, \texttt{not} for $\neg$, \texttt{<->} for
$\leftrightarrow$):
{\footnotesize
\begin{verbatim}
def congrFive (x y : Cyc5.V) : Prop :=
  x.a == y.a [ZMOD 5] /\ x.b == y.b [ZMOD 5] /\
  x.c == y.c [ZMOD 5] /\ x.d == y.d [ZMOD 5]
def unitImage (x : Cyc5.V) : Prop :=
  exists u v : Cyc5.V, Cyc5.mul u v = Cyc5.scalar 1 /\ congrFive x u
def visible (q : Nat) : Prop :=
  q.Prime /\ q % 5 = 1 /\ exists x : Cyc5.V, Cyc5.norm x = (q : Int) /\ not unitImage x
def fullSpectrum : Prop :=
  forall m : Int, (exists f : H -> Int, groupDet f = m) <->
    (m ^ 4 == 1 [ZMOD 125]) \/ (5 ^ 33 : Int) | m \/
      exists n : Int, n != 0 /\ m = 5 ^ 28 * n /\ exists q : Nat, visible q /\ (q : Int) | n
\end{verbatim}}
\noindent It uses \texttt{H} and \texttt{groupDet} from \texttt{Challenge.lean}
and the following definitions from \texttt{Cyc5.lean}, one of the files of the
minimum (verbatim up to transliteration; the docstrings are omitted):
{\footnotesize
\begin{verbatim}
@[ext] structure V where
  a : Int
  b : Int
  c : Int
  d : Int
  deriving DecidableEq
def scalar (n : Int) : V := <n,0,0,0>
def mul (v w : V) : V :=
  let e := v.b*w.d + v.c*w.c + v.d*w.b
  <v.a*w.a+v.c*w.d+v.d*w.c-e, v.b*w.a+v.a*w.b+v.d*w.d-e,
   v.a*w.c+v.b*w.b+v.c*w.a-e, v.a*w.d+v.b*w.c+v.c*w.b+v.d*w.a-e>
def normForm {R : Type*} [CommRing R] (a b c d : R) : R :=
  let s := a*c + b*d + a*d
  let A := a^2 + b^2 + c^2 + d^2 - s
  let B := a*b + b*c + c*d - s
  A^2 - A*B - B^2
def norm (v : V) : Int := normForm v.a v.b v.c v.d
\end{verbatim}}
\noindent An element $\langle a,b,c,d\rangle$ of \texttt{V} stands for
$a+b\zeta+c\zeta^2+d\zeta^3\in R$, in the basis $1,\zeta,\zeta^2,\zeta^3$ of
$R$ over $\Z$; \texttt{mul} is multiplication in $R$ (using
$\zeta^4=-1-\zeta-\zeta^2-\zeta^3$), and \texttt{Cyc5.norm} is the full field
norm $\N$, the product of the four conjugates (proved as
\texttt{Cyc5.norm\_prod}). Because $1,\zeta,\zeta^2,\zeta^3$ is a $\Z$-basis,
\texttt{congrFive} is congruence modulo $5R$; \texttt{unitImage x} says that
the residue of $x$ lies in $U$, a unit being an element with an inverse in
\texttt{V}; and \texttt{visible} is Definition~\ref{def:visible}. The clause
\texttt{q \% 5 = 1} matters: without it, $5$ would count as visible, since
$\N(\pi)=5$ and the residue of $\pi$ is not that of a unit. Like
Definition~\ref{def:visible}, \texttt{visible} uses an existential
quantifier; by Lemma~\ref{lem:choice}, which is proved in the text but not in
Lean, the universal form would define the same primes. \texttt{fullSpectrum} is Theorem~\ref{thm:main}
for every signed integer $m$, with $0$ included in the second set; it uses
\texttt{Cyc5.mul} and \texttt{Cyc5.norm} explicitly, so the ring instances on
\texttt{V} added in later files do not affect its meaning. The final reviewers
restated both frozen statement files with all definitions unfolded (closed by
\texttt{Iff.rfl}); one derived from \texttt{fullSpectrum} in Lean that $11$ is
visible and $5^{28}$ is not a value, so neither side is trivially true.

\paragraph{Review and freezing.} Each statement file was checked against the
source by a Codex agent that had not written it, and frozen by recording its
SHA-256 hash: \texttt{Challenge.lean} (\texttt{9a9a0e18\ldots}) on
21~November 2026 at 01:32~UTC, and \texttt{ChallengeB.lean}
(\texttt{2a5066bf\ldots}) at 03:22~UTC. In both cases the development logs
record no compilation of a proof file before the freeze, and the files were
not changed afterwards; the $37$ files of the minimum were not changed when
the classification was added.

\paragraph{The proofs.} The check files \texttt{Check.lean} and
\texttt{CheckB.lean} contain, apart from \texttt{import} and \texttt{namespace}
lines, exactly
{\footnotesize
\begin{verbatim}
theorem check_witnessValue : Challenge.witnessValue := Proof.witnessValue
theorem check_leastPositiveMultiple : Challenge.leastPositiveMultiple :=
  Proof.leastPositiveMultiple
theorem check_fullSpectrum : ChallengeB.fullSpectrum := ProofB.fullSpectrum
\end{verbatim}}
\noindent together with \texttt{\#print axioms} commands (the line break
inside the second theorem is ours).
The witness value is computed by a certificate: an explicit lower triangular
integer matrix $L$ (\texttt{WitnessCertificate.left}) with
$L\cdot(\text{matrix})$ upper triangular (\texttt{transformed\_eq}, proved
symbolically; \texttt{transformed\_upper}, checked by the kernel), and the two
diagonal products, also evaluated by the kernel, have quotient
$11\cdot5^{28}$ (\texttt{Proof.witnessValue}). Table~\ref{tab:lean} lists the
Lean counterparts of the other steps (names relative to the namespace
\texttt{TaskU.H5}; file names in parentheses, without \texttt{.lean}).

\begin{table}[!htbp]
\centering\footnotesize
\renewcommand{\arraystretch}{1.05}
\begin{tabular}{@{}>{\raggedright\arraybackslash}p{0.40\textwidth}>{\raggedright\arraybackslash}p{0.56\textwidth}@{}}
\hline
Step of the written proof & Lean (file) \\
\hline
Factorization~\eqref{eq:factor} from the matrix &
\texttt{inverseFourier} (\texttt{Fourier}), \texttt{det\_fourier} (\texttt{Blocks}),
\texttt{factorization} (\texttt{Factors}), \texttt{integerFactorization}
(\texttt{IntegerFactors}) \\
Lemma~\ref{lem:normalize} & \texttt{normalize}, \texttt{transport\_det} (\texttt{Normalize}) \\
Lemma~\ref{lem:c0}: $D=c_0+5w$, $\pi\mid w$ &
\texttt{TraceWords.nonlinear\_congruence} (\texttt{TraceCongruence}),
\texttt{TraceWords.integral\_congruence} (\texttt{TraceIntegral}) \\
Proposition~\ref{prop:floor}: $5^{28}\mid M$, the layer data &
\texttt{aug\_necessary}, \texttt{layers\_exist}, \texttt{Layers.value},
\texttt{Layers.deep\_zero}, \texttt{Layers.deep\_nonlin} (\texttt{Divisibility});
\texttt{Proof.baseline} (\texttt{Minimum}) \\
Fact~(d), Lemma~\ref{lem:phase} &
\texttt{Cyc5.unit\_phase}, \texttt{Cyc5.ordinary\_phase} (\texttt{CycPhase});
\texttt{phases\_vanish} (\texttt{RayMoments}); \texttt{BTrace.exact\_trace\_zero} (\texttt{BTrace}) \\
Corollary~\ref{cor:min}, lower bound: for $m\le10$ the inert primes
$2,3,7$ force the $q_j$ to be units &
\texttt{Cyc5.inert\_fourth} (\texttt{CycNorm}); \texttt{Layers.small\_norm} (\texttt{Divisibility});
\texttt{ordinary\_trace\_zero} (\texttt{TraceLower});
\texttt{TraceFinal.deep\_average} (\texttt{TraceFinal}); \texttt{Proof.small},
\texttt{Proof.leastPositiveMultiple} (\texttt{Minimum});
\texttt{SmallMultiples.least\_from\_small} (\texttt{SmallMultiples}) \\
Lemma~\ref{lem:coprime} &
\texttt{Coprime.determinant\_congruence}, \texttt{Coprime.coprime\_necessary}
(\texttt{CoprimeCongruence}) \\
Fact~(c) & instance \texttt{IsPrincipalIdealRing V} (\texttt{PrimeRing}), from
Mathlib's \texttt{IsCyclotomicExtension.Rat.five\_pid} \\
Lemma~\ref{lem:units} &
\texttt{Spectrum.inv\_table}, \texttt{Spectrum.normal\_table}, \texttt{Spectrum.fixed\_table}
(\texttt{UnitTable}); \texttt{Spectrum.conjugation\_fixed},
\texttt{Spectrum.visible\_normalized} (\texttt{UnitResidues}) \\
Lemma~\ref{lem:extract} &
\texttt{PrimeExtract.prime\_visible}, \texttt{PrimeExtract.visible\_dvd\_of\_not\_unitImage}
(\texttt{PrimeExtract}); \texttt{PrimeOrbit.four\_dvd} (\texttt{PrimeOrbit}) \\
Replacement step of Proposition~\ref{prop:necessity} &
\texttt{BTrace.trace\_replacement} (\texttt{BTrace}), \texttt{Spectrum.deep\_of\_images}
(\texttt{BTraceImage}) \\
Proposition~\ref{prop:necessity} & \texttt{Spectrum.divisible\_necessary} (\texttt{Necessity}) \\
Seeds and~\eqref{eq:rows} &
\texttt{Spectrum.seedData}, \texttt{Spectrum.seed\_residues} (\texttt{Seeds});
\texttt{Spectrum.swapped\_data} (\texttt{SeedMarginals}); \texttt{Spectrum.rows\_value} (\texttt{RowsValue}) \\
Lemma~\ref{lem:lift} &
\texttt{Spectrum.lift\_shift} (\texttt{LiftShift}); \texttt{Spectrum.compact\_matrix},
\texttt{Spectrum.baseBlock\_det}, \texttt{Spectrum.minorDir}, \texttt{Spectrum.minorResidue}
(\texttt{LiftColumns}); \texttt{Spectrum.shift\_block} (\texttt{LiftCorrection}) \\
Propositions~\ref{prop:tail} and~\ref{prop:low} &
\texttt{Spectrum.insert\_rows}, \texttt{Spectrum.seed\_quot\_spec} (\texttt{Insert});
\texttt{Spectrum.high\_sufficient}, \texttt{Spectrum.low\_sufficient} (\texttt{Construct}) \\
Proposition~\ref{prop:coprime} &
\texttt{Coprime.seed\_det}, \texttt{Coprime.det\_neg}, \texttt{Coprime.coprime\_sufficient}
(\texttt{CoprimeSeeds}) \\
Theorem~\ref{thm:main} & \texttt{ProofB.fullSpectrum} (\texttt{ProofB}) \\
\hline
\end{tabular}
\caption{Lean counterparts of the steps of the proofs.}
\label{tab:lean}
\end{table}

\noindent The Lean proofs do not assume the results quoted
from~\cite{mp} and~\cite{pan25}: the factorization, the congruence of
Lemma~\ref{lem:c0}, the bound $5^{28}$, both directions of~\cite[Corollary~3.2]{mp}
and the ray norms of the seeds are proved again (the seed values by kernel
computation from the explicit coefficients). The class number of $K$ is taken
from Mathlib's proved theorem. The axioms of \texttt{Check.check\_witnessValue},
\texttt{Check.check\_leastPositiveMultiple}, \texttt{check\_fullSpectrum} and
the corresponding \texttt{Proof} and \texttt{ProofB} theorems are exactly
\texttt{propext}, \texttt{Classical.choice} and \texttt{Quot.sound}. No file uses
\texttt{sorry}, \texttt{native\_\allowbreak decide},
\texttt{bv\_\allowbreak decide} or an \texttt{axiom} declaration; all uses of
\texttt{decide}, including the large tables, are checked by the kernel.

\paragraph{Checks.} The $37$ files of the minimum were compiled module by
module in two new empty directories on the author's Mac, with the pinned
toolchain and with warnings treated as errors, on 21~November 2026
(01:42:55--01:56:09~UTC); each directory also ran one build of the whole
package, and every module was replayed with \texttt{leanchecker}. All $150$
commands ($75$ per directory) ended with exit code~$0$, the
\texttt{leanchecker} runs printed nothing, and the two directories produced
identical \texttt{.olean} files. A first attempt, started $35$ seconds earlier, was
stopped by the acceptance script, which misread two informational messages of
Lake; it was kept on record, and the script was corrected only in this
respect. The $63$ files of the classification ($37+26$) were checked in the
same way in two further new directories (03:33:17--03:55:22~UTC): $254$
commands ($127$ per directory), all with exit code~$0$, empty
\texttt{leanchecker} output and identical \texttt{.olean} files. The final
reviewers (see the next section) compiled the files again with plain
\texttt{lean} in new directories of their own, together with short audit files
of their own, two of them with warnings treated as errors and the third
reporting no warnings, and replayed every module with
\texttt{leanchecker}: one reviewer the $37$ files of the minimum (03:21--03:27~UTC),
and two reviewers independently the $63$ files of the classification
(04:29--04:41 and 04:42--04:55~UTC); all runs on the project files ended with
exit code~$0$ (one line of a reviewer's own audit file had to be changed and
recompiled). Mathlib
and its dependencies were taken from a shared prebuilt cache at the pinned
commits and were not rechecked; the \texttt{leanchecker} replays cover the
project modules. In the repository the files have module names under
\texttt{Research.TaskU.H5} and \texttt{Research.TaskU.H5B} and differ from the
checked files only in their \texttt{import} lines (the four files that
import only Mathlib, among them \texttt{Challenge.lean}, are unchanged).
The repository versions were made by a script that changed only the local \texttt{import} lines, and on 21 November 2026 (17:02--17:22~UTC) the $63$ modules were compiled and replayed with \texttt{leanchecker} again, module by module, in two new empty directories on the author's Mac, one after the other; there were no warnings, and the two directories produced identical \texttt{.olean} files and the same axiom lists.
In the whole-project rebuild of version~\repoversion{} of the repository (on the author's Windows computer, in two passes on 20 and 21 November 2026), these files were compiled and replayed with \texttt{leanchecker} again in pass~2 (21 November 2026), which built the files not included in pass~1 on top of its results; no module needed a second attempt (details in the repository). The formalization has not been checked with comparator or with
an independent kernel.

\par}

\section*{How this paper was produced}
This work was carried out on 21 and 22 November 2026 (UTC) in a workflow directed by
the author, using OpenAI Codex with the model GPT-6.2-Astra (reasoning level
ultra) and Anthropic Claude Code with the model Claude Opus~5.5. The author
asked for a new search for substantial open problems, and Claude Code passed
this request to Codex at 01:00~UTC. In that search Codex screened $72$
problems not screened in earlier rounds and prepared notes on eight of them,
this one among them. According to the file times, by 01:12~UTC a Codex agent
had written the proof of the minimum (Corollary~\ref{cor:min}) and checked the
value of the witness by exact elimination on the full matrix, and another
Codex agent (same model) that had not derived it had read the proof,
recomputed the witness value by the formulas of~\cite{mp} and reran the
elimination program,
and approved it; the search ended at 01:22~UTC. The author chose this problem
and one other. Shortly before 01:28~UTC, Codex agents (same model) began to
check the reading of the source, search again for prior work, write the
statement files, the Lean proofs and the written proofs, and run the checks;
the review and freezing of the statements are described in
Section~\ref{sec:lean}. According to the file timestamps and the development
logs, the proof files of the minimum were created after its freeze and
compiled during development between 01:32 and 01:42~UTC. In parallel, Codex
wrote the written classification (file times 01:36 to 01:41~UTC), which was
read by Codex agents that had not derived it (reviews at 01:39, 01:42 and
01:44~UTC); one of them rebuilt the coefficients of $25$ constructed elements and recomputed
their full $125\times125$ determinants. As allowed by the task description,
Claude Code asked Codex at 03:20~UTC to formalize the classification; its proof
files were created after the freeze of \texttt{ChallengeB.lean} and compiled
during development between 03:23 and 03:33~UTC. The Lean route follows the
written classification, except that Lemma~\ref{lem:lift} uses explicit
cofactor certificates in place of a general lifting argument.

The final reviews were done by Claude Code agents (model Claude Opus~5.5) that
had written none of the files: one for the minimum (03:20 to about 03:28~UTC)
and, because the classification is the main result, two independent ones for
it (03:57 to 04:55~UTC), neither of which read the other's report. Each
reviewer read the source itself, compared the statement files with it, read
the written proofs and the proof files, checked the receipts of the checks,
and rebuilt the formalization as described above. The two reviewers of the classification tested it on random integer
functions $f$ ($4{,}650$ exact determinants in one review; about $300{,}000$
evaluations through~\eqref{eq:factor}, checked against $140$ full
determinants, in the other; about $150{,}000$ cases with valuation $28$ to
$32$) and found no counterexample. All
three decisions were APPROVE\_FORMAL, with no blocking or major remarks; their
minor remarks ($6$, $5$ and $5$, some of them optional) concern the written
documents; those that apply to this paper are taken into account here, among them the naming of the Lean definitions, the
basis and the norm, the choice independence of visibility, the status of
Corollary~\ref{cor:33}, the citation of~\cite{mp}, its example for $p=5$, and
the labelling of Remark~\ref{rem:cft}. Claude Code (Claude Opus~5.5) checked
the sources, recomputed in Python the determinants of the witness and of the
example of~\cite{mp}, the group $U$ and the visible primes below $6000$, and
wrote this text from Codex's manuscripts and the Lean files; the class field
reformulation in Remark~\ref{rem:cft}, the proof sketch of
Lemma~\ref{lem:coprime} and the arrangement of the necessity proof around
Lemma~\ref{lem:extract} were added while writing and were not part of the
documents read by the final reviewers. An independent Claude Code agent (model Claude Opus~5.5) that had written none of the files reviewed this paper, and another reviewed the release; their corrections were applied. Throughout, ``we''
refers to this workflow; the reading and searching described in the paper
were done by the agents. Codex agents searched for prior work in the search for
problems, again before the statement of the minimum was frozen (01:28--01:29~UTC,
with further searches until 01:42~UTC), and for the classification in one
query just before and in the main search shortly after its statement was frozen
(03:21 and 03:28--03:31~UTC); the final reviewers
searched again; up to 21~November 2026 these searches found no earlier
determination of the multiples of $5$ in $\cS(H_5)$, of its least positive
multiple of $5$, or of a prime-factor criterion of the kind in
Theorem~\ref{thm:main} for $H_5$ (for $C_7\times C_7$ see~\cite{pan49}; details
in the repository). The journal version of~\cite{mp} could not be
read, and one citation index refused some queries (rate limit), so a further
search, including recent Zenodo records of group determinant classifications,
was planned for the release. Claude Code made this search on 21 November 2026 (17:54--18:21~UTC; the arXiv, Semantic Scholar, scite, Crossref, DataCite, Zenodo, including the series of classifications of group determinants of small orders by A.~Radolovi\'c, GitHub and repositories of AI-assisted mathematics) and found no new related work; the same citation index again refused to list the works citing~\cite{mp} (rate limit). Claude Code searched again on 22 November 2026 (05:39--05:51~UTC), when the OpenAlex lists of citing works refused on 21 November were retrieved, and found nothing new. This reports the
coverage of the searches, not a guarantee that they were complete. The author
is not a professional mathematician. The author has read this paper and, using a side-by-side table, compared its Lean statements with the main theorems: Theorem~\ref{thm:main} and Corollary~\ref{cor:min}. No human mathematician has
read the proofs or reviewed this paper. The correctness of
Theorem~\ref{thm:main} and Corollary~\ref{cor:min} rests on the Lean proofs and
on the agreement of the Lean statements with the definitions of~\cite{mp}; the
written proofs are given for the reader. AI systems are not authors of this
paper; the author takes full responsibility for its content.

\medskip
\noindent Thanks are due to Michael J. Mossinghoff and Christopher Pinner for
the problem and the methods on which this work rests, and to Chatchawan
Panraksa for the seeds.

{\small
\begin{thebibliography}{9}
\setlength{\itemsep}{1pt}
\bibitem{bhj} W.~Bley, T.~Hofmann and H.~Johnston, Computation of lattice
isomorphisms and the integral matrix similarity problem, \emph{Forum Math.
Sigma} 10 (2022), e87, \url{https://doi.org/10.1017/fms.2022.74};
arXiv:2202.03526.
\bibitem{repo} H.~Dong, \emph{ai4math-results}, repository,
\url{https://github.com/Horace-Maxwell/ai4math-results}, version~\repoversion{}
(2026); archived at Zenodo under the concept DOI
\url{https://doi.org/10.5281/zenodo.22970254}.
\bibitem{mathlib}
The mathlib Community, The Lean mathematical library, in \emph{Proc.\ 9th ACM
SIGPLAN Int.\ Conf.\ Certified Programs and Proofs} (CPP '20), ACM, 2020,
\url{https://doi.org/10.1145/3372885.3373824}.
\bibitem{mp} M.~J. Mossinghoff and C.~Pinner, The integer group determinants
for the Heisenberg group of order $p^3$, \emph{Michigan Math.\ J.} 74 (2024),
no.~3, 551--569, \url{https://doi.org/10.1307/mmj/20216124};
arXiv:2108.04870v1 (2021).
\bibitem{pan25} C.~Panraksa, The $5$-divisible integer group determinants for
the elementary abelian group of order $25$, arXiv:2608.04267v1 (2026).
\bibitem{pan49} C.~Panraksa, Integer group determinants for the elementary
abelian group of order $49$, arXiv:2609.09542v1 (2026).
\bibitem{panp} C.~Panraksa, Bernoulli numbers and the integer group
determinants of $C_p\times C_p$, arXiv:2610.15262v1 (2026).
\bibitem{ps} C.~Pinner and C.~Smyth, Integer group determinants for small
groups, \emph{Ramanujan J.} 51 (2020), no.~2, 421--453,
\url{https://doi.org/10.1007/s11139-018-0092-x}; arXiv:1806.00199v1.
\end{thebibliography}}
\end{document}
